% ================================================================
%  conteudo/conclusao.tex
% ================================================================

\section{Conclusão}

A conclusão é a seção final do relatório e deve retomar, de forma sintética, os principais achados da prática experimental, respondendo diretamente aos objetivos propostos. Ela não deve apresentar dados novos nem repetir extensamente a discussão — seu papel é consolidar o que foi aprendido e indicar o alcance e as limitações dos resultados obtidos.

Com base nos resultados obtidos, foi possível verificar que a concentração inicial de substrato exerceu influência significativa sobre a cinética de crescimento de \textit{Saccharomyces cerevisiae} nas condições experimentais avaliadas. O aumento da concentração de glicose proporcionou incremento na produção de biomassa, porém sem aumento proporcional da velocidade específica máxima de crescimento ($\mu_{max}$), indicando possível efeito de inibição por substrato em concentrações mais elevadas.

Os parâmetros cinéticos determinados a partir das curvas de crescimento demonstraram boa concordância com os valores reportados na literatura para a mesma cepa microbiana, validando a metodologia empregada. O método do DNS mostrou-se adequado para a quantificação de glicose residual nas faixas de concentração estudadas, com coeficiente de variação inferior a 5\% entre as triplicatas analisadas.

Conclui-se que a prática atingiu os objetivos propostos, permitindo a compreensão prática dos conceitos de cinética microbiana abordados na disciplina. Sugere-se, para trabalhos futuros, a avaliação de outras fontes de carbono e a realização do processo em modo contínuo, visando à comparação dos parâmetros cinéticos obtidos nos dois modos de operação.

